% GENERATED FILE. Edit canonical lessons, then run npm run generate:books.
% canonical-source-hash: fnv1a64:851040a9cea3b3da
% canonical-lessons: GU-C322-humlo, GU-C322-ladai, GU-C322-yoddho, GU-C322-hanvu, GU-C322-jhajhumvu

\chapter{The Attack at Dawn}
\label{ch:gu-the-attack-at-dawn}
\hlchaptermodality{322}

\begin{chapteropening}
\textbf{By the end of this chapter:} \emph{I can tell how an attack began and how long the fight went on.}

Complete the last lesson of chapter 322: I can tell how an attack began and how long the fight went on.
\end{chapteropening}

\section[\texorpdfstring{\gu{હુમલો} (humlo)}{હુમલો (humlo)}]{\gu{હુમલો} (humlo) — an attack}

\label{lesson:GU-C322-humlo}



\begin{quote}
Before the new one: say the Gujarati for a cannon, then the Gujarati for to pierce.
\end{quote}

\subsection*{You'll want to know: \gu{હુમલો}}
\textbf{\gu{હુમલો}} — \emph{humlo} — \textquotedblleft{}an attack\textquotedblright{}.

\begin{quote}
\gu{પરોઢે} \gu{દુશ્મને} \gu{કિલ્લા} \gu{પર} \gu{હુમલો} \gu{કર્યો}.
\end{quote}

\begin{quote}
\emph{paroḍhe duśmane killā par humlo karyo.}
\end{quote}

\begin{quote}
At dawn the enemy attacked the fort.
\end{quote}

A masculine noun.

\begin{cousinweb}
\textbf{\gu{હુમલો}} (\emph{humlo}) is usually traced, through Persian, to Arabic \emph{ḥamla}, \textquotedblleft{}a charge, an attack\textquotedblright{}. The place attacked takes \textbf{\gu{પર}}: \textbf{\gu{કિલ્લા} \gu{પર} \gu{હુમલો}}.
\end{cousinweb}

\begin{grammarlens}[title={What you've built}]
One more word for telling what happened.
\end{grammarlens}

\subsection*{Your turn}
\begin{itemize}
\raggedright
  \item \emph{Say it:} \emph{humlo}
  \item \emph{Say it:} \emph{humlo}, once more
  \item \emph{Say it:} \emph{paroḍhe duśmane killā par humlo karyo.} — At dawn the enemy attacked the fort.
  \item \emph{Recall:} say the Gujarati for a cannon, then the Gujarati for to pierce, then say \emph{humlo} again
\end{itemize}

\begin{tcolorbox}[breakable,colback=teal!4,colframe=teal!35!black,title={Before you move on}]
What does \gu{હુમલો} mean? (\textquotedblleft{}An attack\textquotedblright{}.) Say it once more.
\end{tcolorbox}
\section[\texorpdfstring{\gu{લડાઈ} (laḍāī)}{લડાઈ (laḍāī)}]{\gu{લડાઈ} (laḍāī) — a battle; a fight}

\label{lesson:GU-C322-ladai}



\begin{quote}
Before the new one: say the Gujarati for to pierce, then the Gujarati for an attack.
\end{quote}

\subsection*{You'll want to know: \gu{લડાઈ}}
\textbf{\gu{લડાઈ}} — \emph{laḍāī} — \textquotedblleft{}a battle; a fight\textquotedblright{}.

\begin{quote}
\gu{આ} \gu{લડાઈ} \gu{આખો} \gu{દિવસ} \gu{ચાલી}.
\end{quote}

\begin{quote}
\emph{ā laḍāī ākho divas cālī.}
\end{quote}

\begin{quote}
This battle went on all day.
\end{quote}

A feminine noun.

\begin{cousinweb}
\textbf{\gu{લડાઈ}} (\emph{laḍāī}) is built on \textbf{\gu{લડવું}}, to fight, from \emph{*laḍ}, a form CDIAL reconstructs (CDIAL 10920), with the noun ending \emph{-āī}: the fighting itself, from a quarrel to a battle.
\end{cousinweb}

\begin{grammarlens}[title={What you've built}]
One more word for telling what happened.
\end{grammarlens}

\subsection*{Your turn}
\begin{itemize}
\raggedright
  \item \emph{Say it:} \emph{laḍāī}
  \item \emph{Say it:} \emph{laḍāī}, once more
  \item \emph{Say it:} \emph{ā laḍāī ākho divas cālī.} — This battle went on all day.
  \item \emph{Recall:} say the Gujarati for to pierce, then the Gujarati for an attack, then say \emph{laḍāī} again
\end{itemize}

\begin{tcolorbox}[breakable,colback=teal!4,colframe=teal!35!black,title={Before you move on}]
What does \gu{લડાઈ} mean? (\textquotedblleft{}A battle; a fight\textquotedblright{}.) Say it once more.
\end{tcolorbox}
\section[\texorpdfstring{\gu{યોદ્ધો} (yoddho)}{યોદ્ધો (yoddho)}]{\gu{યોદ્ધો} (yoddho) — a warrior}

\label{lesson:GU-C322-yoddho}



\begin{quote}
Before the new one: say the Gujarati for an attack, then the Gujarati for a battle.
\end{quote}

\subsection*{You'll want to know: \gu{યોદ્ધો}}
\textbf{\gu{યોદ્ધો}} — \emph{yoddho} — \textquotedblleft{}a warrior\textquotedblright{}.

\begin{quote}
\gu{ગામના} \gu{યોદ્ધા} \gu{બહાદુર} \gu{હતા}.
\end{quote}

\begin{quote}
\emph{gāmnā yoddhā bahādur hatā.}
\end{quote}

\begin{quote}
The village's warriors were brave.
\end{quote}

A masculine noun.

\begin{cousinweb}
goes back to Sanskrit \emph{yoddhṛ}, \textquotedblleft{}a fighter, a warrior\textquotedblright{} (Monier-Williams), made into a Gujarati noun in \emph{-o}: \textbf{\gu{યોદ્ધો}}, plural \textbf{\gu{યોદ્ધા}}.
\end{cousinweb}

\begin{grammarlens}[title={What you've built}]
One more word for telling what happened.
\end{grammarlens}

\subsection*{Your turn}
\begin{itemize}
\raggedright
  \item \emph{Say it:} \emph{yoddho}
  \item \emph{Say it:} \emph{yoddho}, once more
  \item \emph{Say it:} \emph{gāmnā yoddhā bahādur hatā.} — The village's warriors were brave.
  \item \emph{Recall:} say the Gujarati for an attack, then the Gujarati for a battle, then say \emph{yoddho} again
\end{itemize}

\begin{tcolorbox}[breakable,colback=teal!4,colframe=teal!35!black,title={Before you move on}]
What does \gu{યોદ્ધો} mean? (\textquotedblleft{}A warrior\textquotedblright{}.) Say it once more.
\end{tcolorbox}
\section[\texorpdfstring{\gu{હણવું} (haṇvũ)}{હણવું (haṇvũ)}]{\gu{હણવું} (haṇvũ) — to slay}

\label{lesson:GU-C322-hanvu}



\begin{quote}
Before the new one: say the Gujarati for a battle, then the Gujarati for a warrior.
\end{quote}

\subsection*{You'll want to know: \gu{હણવું}}
\textbf{\gu{હણવું}} — \emph{haṇvũ} — \textquotedblleft{}to slay\textquotedblright{}.

\begin{quote}
\gu{યોદ્ધાએ} \gu{એક} \gu{વાઘને} \gu{હણ્યો}.
\end{quote}

\begin{quote}
\emph{yoddhāe ek vāghne haṇyo.}
\end{quote}

\begin{quote}
The warrior slew a tiger.
\end{quote}

Past, with \gu{તેણે} (he or she …): \gu{તેણે} \gu{દુશ્મનને} \gu{હણ્યો}, \emph{teṇe duśmanne haṇyo}.

\begin{cousinweb}
\textbf{\gu{હણવું}} (\emph{haṇvũ}) is inherited from Sanskrit \emph{hánati}, \textquotedblleft{}strikes, kills\textquotedblright{}, and CDIAL lists the Gujarati verb, \textquotedblleft{}to strike, kill\textquotedblright{} (CDIAL 13963). It is a word of old stories; \textbf{\gu{મારવું}} is the everyday verb.
\end{cousinweb}

\begin{grammarlens}[title={What you've built}]
One more word for telling what happened.
\end{grammarlens}

\subsection*{Your turn}
\begin{itemize}
\raggedright
  \item \emph{Say it:} \emph{haṇvũ}
  \item \emph{Say it:} \emph{haṇvũ}, once more
  \item \emph{Say it:} \emph{yoddhāe ek vāghne haṇyo.} — The warrior slew a tiger.
  \item \emph{Recall:} say the Gujarati for a battle, then the Gujarati for a warrior, then say \emph{haṇvũ} again
\end{itemize}

\begin{tcolorbox}[breakable,colback=teal!4,colframe=teal!35!black,title={Before you move on}]
What does \gu{હણવું} mean? (\textquotedblleft{}To slay\textquotedblright{}.) Say it once more.
\end{tcolorbox}
\section[\texorpdfstring{\gu{ઝઝૂમવું} (jhajhūmvũ)}{ઝઝૂમવું (jhajhūmvũ)}]{\gu{ઝઝૂમવું} (jhajhūmvũ) — to fight on; to struggle}

\label{lesson:GU-C322-jhajhumvu}



\begin{quote}
Before the new one: say the Gujarati for a warrior, then the Gujarati for to slay.
\end{quote}

\subsection*{You'll want to know: \gu{ઝઝૂમવું}}
\textbf{\gu{ઝઝૂમવું}} — \emph{jhajhūmvũ} — \textquotedblleft{}to fight on; to struggle\textquotedblright{}.

\begin{quote}
\gu{સૈનિકો} \gu{છેલ્લી} \gu{પળ} \gu{સુધી} \gu{ઝઝૂમ્યા}.
\end{quote}

\begin{quote}
\emph{sainiko chellī paḷ sudhī jhajhūmyā.}
\end{quote}

\begin{quote}
The soldiers fought on to the very last moment.
\end{quote}

Past (he …): \gu{તે} \gu{ખૂબ} \gu{ઝઝૂમ્યો}, \emph{te khūb jhajhūmyo}.

\begin{cousinweb}
\textbf{\gu{ઝઝૂમવું}} (\emph{jhajhūmvũ}) probably belongs with \emph{jhujhvũ}, to fight, which CDIAL lists for Gujarati under Sanskrit \emph{yúdhyatē}, \textquotedblleft{}fights against\textquotedblright{} (CDIAL 10502); CDIAL does not list this longer form, so the link is not settled. It is to keep fighting against long odds: \textbf{\gu{છેલ્લી} \gu{પળ} \gu{સુધી} \gu{ઝઝૂમ્યા}}.
\end{cousinweb}

\begin{grammarlens}[title={What you've built}]
One more word for telling what happened.
\end{grammarlens}

\subsection*{Your turn}
\begin{itemize}
\raggedright
  \item \emph{Say it:} \emph{jhajhūmvũ}
  \item \emph{Say it:} \emph{jhajhūmvũ}, once more
  \item \emph{Say it:} \emph{sainiko chellī paḷ sudhī jhajhūmyā.} — The soldiers fought on to the very last moment.
  \item \emph{Recall:} say the Gujarati for a warrior, then the Gujarati for to slay, then say \emph{jhajhūmvũ} again
\end{itemize}

\begin{tcolorbox}[breakable,colback=teal!4,colframe=teal!35!black,title={Before you move on}]
What does \gu{ઝઝૂમવું} mean? (\textquotedblleft{}To fight on; to struggle\textquotedblright{}.) Say it once more.
\end{tcolorbox}
